\section{Problem formulation}\label{sec:problem}

This section formalises selective validation as a three-way routing problem. It states the guarantee that the
accept rule provides and the assumption on which that guarantee rests. It then derives the exact cost model
used to judge whether screening pays off.

\subsection{Agent, validator and episode}
An \emph{episode} begins when a monitoring layer detects that the process has left its safe operating region,
or, in path planning, when a query is posed. In round $k=0,1,\dots,K-1$ the agent receives a prompt $x_k$ and
returns a proposal $a_k$. The first prompt $x_0$ describes the task state. A later prompt $x_k$ ($k\ge 1$)
additionally contains feedback on the previous proposal. The \emph{validator} $V(x,a)\in\{0,1\}$ returns
whether proposal $a$ is acceptable. In the CSTR it simulates the plant from the current state under the proposed
setpoints and checks the outcome contract (Section~\ref{sec:setup}). We write $Y=1-V$ for the failure
indicator. In the conventional loop, which we call \emph{always-validate} (AV), every proposal is validated. A
passing proposal is executed and ends the episode. A failing one is returned to the agent with the validator's
measured outcome as feedback, until a budget of $K=6$ proposals is exhausted. If no proposal passes, the current
setpoints are held.

While generating $a_k$, the agent's forward pass produces an internal record $h_k$: output-token
distributions, attention weights and hidden states. A \emph{signal map} $\psi$ turns the available information
into a feature vector $z=\psi(x,a,h)$. The signal families are defined in Section~\ref{sec:signals}. Some use
only $h$ (internal signals), some use only $(x,a)$ (observable signals), and some use both.

\subsection{Three-way routing}
A \emph{screen} consists of a probe $f$, which maps $z$ to an estimated failure probability $\hat r=f(z)\in[0,1]$,
and two thresholds $\tacc<\trej$:
\begin{equation}\label{eq:route}
d(\hat r)=\begin{cases}
\ACC & \hat r\le \tacc \quad\text{(execute without validation)},\\
\REJ & \hat r\ge \trej \quad\text{(return to the agent without validation)},\\
\VAL & \text{otherwise}.
\end{cases}
\end{equation}
Both \ACC{} and \REJ{} avoid a validator call. The two errors differ in kind. An \ACC{} of a failing proposal
executes an unsafe action on the plant; this is the safety-relevant error. A \REJ{} of a passing proposal costs
a regeneration and may cost the episode its recovery if the budget runs out; this is an efficiency error. The
design therefore treats the two thresholds asymmetrically: $\tacc$ carries a statistical guarantee, while
$\trej$ is set for high precision.

\begin{definition}[Offline operating characteristics]\label{def:metrics}
For a screen applied to proposals drawn from a distribution $P$, the \emph{skip rate} is
$S=P\{d\in\{\ACC,\REJ\}\}$, the \emph{unchecked failure rate} is $\rho_{A}=P\{Y=1\mid d=\ACC\}$, and the
\emph{rejection precision} is $\pi_{R}=P\{Y=1\mid d=\REJ\}$.
\end{definition}

The design goal is to maximise $S$ subject to $\rho_A\le X$ with high confidence, for a tolerance $X$ set by the
operator, and subject to a high rejection precision. We use $X=10\%$ as the primary tolerance and $X=5\%$ as a
secondary one. The validator calls skipped, $S$, are the primary endpoint of the study. We use this rather than
the area under the ROC curve (AUROC) because $S$ measures what an engineer saves at a fixed risk, whereas AUROC
averages over operating points that would never be deployed.

\subsection{Choosing the thresholds}
Let $\mathcal D=\{(\hat r_i,Y_i)\}$ be development proposals that were not used to fit the probe. For a
candidate threshold $t$ let $m(t)=|\{i:\hat r_i\le t\}|$ be the number of proposals it would accept and $k(t)$ the
number of failures among them. Let
\begin{equation}\label{eq:cp}
U_{\alpha}(k,m)=\begin{cases}\mathrm{Beta}^{-1}(1-\alpha;\;k+1,\;m-k) & k<m,\\ 1 & k=m,\end{cases}
\end{equation}
be the one-sided Clopper--Pearson upper confidence bound at level $1-\alpha$ on a binomial proportion with $k$
successes in $m$ trials \citep{clopper1934}; $U_\alpha(0,0)=1$. The accept threshold is the largest threshold
whose bound does not exceed the tolerance,
\begin{equation}\label{eq:tacc}
\tacc=\max\{t:\;U_{\delta_{\mathrm{sel}}}(k(t),m(t))\le X\},\qquad \delta_{\mathrm{sel}}=0.05,
\end{equation}
and the reject threshold is the smallest threshold whose rejected development set is at least 95\% failures,
\begin{equation}\label{eq:trej}
\trej=\min\Bigl\{t:\;\tfrac{1}{|\{i:\hat r_i\ge t\}|}\textstyle\sum_{i:\hat r_i\ge t}Y_i\ge 0.95\Bigr\}.
\end{equation}
If no threshold satisfies Eq.~\eqref{eq:tacc}, the screen accepts nothing. Fig.~\ref{fig:threshold} illustrates
why the bound matters. A rule that sets $\tacc$ where the \emph{point estimate} $k(t)/m(t)$ reaches $X$ selects a
larger threshold. The failure rate of the set it selects is then biased upward on new data, because the selection
itself favours thresholds whose development failure rate happened to be low. In an earlier pre-registered stage
of this study the point rule did prove too loose on certification data. The bound-based rule replaced it by a
dated amendment, made before any affected test data were analysed.

\begin{figure}[!htbp]
\centering
% Schematic of the bound-based accept rule (curves are illustrative, not data).
\begin{tikzpicture}
\begin{axis}[paperaxis, width=0.62\textwidth, height=0.36\textwidth,
  xmin=0, xmax=0.62, ymin=0, ymax=20, domain=0.03:0.6, samples=80,
  xlabel={accept threshold $t$ on the predicted risk $\hat r$},
  ylabel={failure rate of $\{\hat r\le t\}$ (\%)},
  xtick={0,0.2,0.4,0.6}, ytick={0,5,10,15,20}, ymajorgrids=false,
  legend pos=north west, legend cell align=left, clip=false]
\addplot[cObs, thick] {30*x^1.5};
\addlegendentry{point estimate on development data}
\addplot[cInt, thick, densely dashed] {30*x^1.5 + 2.5/sqrt(x)};
\addlegendentry{one-sided Clopper--Pearson upper bound}
\addplot[black!70, thin, domain=0:0.62] {10};
\node[anchor=north east, font=\scriptsize] at (axis cs:0.61,10) {tolerance $X$};
% chosen thresholds
\draw[cInt, thin] (axis cs:0.326,0) -- (axis cs:0.326,10);
\fill[cInt] (axis cs:0.326,10) circle (1.6pt);
\node[anchor=north west, font=\scriptsize, text=cInt, align=left] at (axis cs:0.33,3.6) {$t_{\mathrm{acc}}$\\[-1pt](bound rule)};
\draw[cObs, thin] (axis cs:0.481,0) -- (axis cs:0.481,10);
\fill[cObs] (axis cs:0.481,10) circle (1.6pt);
\node[anchor=north west, font=\scriptsize, text=cObs, align=left] at (axis cs:0.485,3.6) {point rule\\[-1pt](too loose)};
\end{axis}
\end{tikzpicture}

\caption{Schematic of the bound-based accept rule (illustrative curves, not data). For each candidate threshold
$t$, the failure rate of the development proposals with $\hat r\le t$ (solid) and its one-sided Clopper--Pearson
upper bound (dashed) are computed. The bound is large for small $t$, because few proposals are accepted, and
approaches the point estimate as more are accepted. The accept threshold $\tacc$ is the largest $t$ at which the
bound stays below the tolerance $X$. The point-estimate rule picks a larger threshold and, after selection,
exceeds $X$ on new data.}
\label{fig:threshold}
\end{figure}

\subsection{What the certificate guarantees}
After $f$, $\tacc$ and $\trej$ are frozen, each screen is evaluated once on a certification set $\mathcal C$ that
played no role in building it. Let $m$ be the number of certification proposals it accepts and $k$ the number
of failures among them.

\begin{proposition}[Certified unchecked failure rate]\label{prop:cert}
Let the screen be fixed independently of $\mathcal C$, and let the proposals in $\mathcal C$ and the future
proposals on which the screen is deployed be independent draws from the same distribution $P$. Then, for the
deployed unchecked failure rate $\rho_A=P\{Y=1\mid \hat r\le\tacc\}$,
\[
\Pr\bigl\{\rho_A\le U_{\delta/M}(k,m)\bigr\}\ \ge\ 1-\delta/M,
\]
and the statements for $M$ screens certified on the same set hold simultaneously with probability at least
$1-\delta$.
\end{proposition}
\begin{proof}
Because the screen does not depend on $\mathcal C$, the accept decisions are a fixed function of each
certification proposal's features. Conditional on the set of accepted proposals, their failure indicators are
therefore independent Bernoulli variables with mean $\rho_A$, so $k\mid m\sim\mathrm{Bin}(m,\rho_A)$. The
Clopper--Pearson bound covers a binomial proportion with probability at least $1-\delta/M$ for every $m$
\citep{clopper1934}. The same coverage therefore holds unconditionally. The simultaneous statement follows from
the union (Bonferroni) bound.
\end{proof}

We call a screen \emph{certified} at tolerance $X$ if $U_{\delta/M}(k,m)\le X$ with $\delta=0.05$ and $M$ the number
of signal sets compared on that certification set ($M=10$ in path planning, $M=11$ in the CSTR).
Proposition~\ref{prop:cert} has two practical consequences.

\emph{A minimum accept-set size.} With no failures among the accepts, $U_\alpha(0,m)=1-\alpha^{1/m}$, so
certification at tolerance $X$ requires
\begin{equation}\label{eq:mmin}
m\ \ge\ m_{\min}=\left\lceil\frac{\ln\alpha}{\ln(1-X)}\right\rceil .
\end{equation}
For $X=10\%$ and $\alpha=0.05/11$, $m_{\min}=52$; any failure among the accepts raises the requirement. A
certification set of 500 episodes in which a screen accepts 10\% of proposals therefore cannot certify it, however
safe the accepts are. We stated this power limit for the CSTR in advance.

\emph{An exchangeability premise.} The guarantee covers proposals drawn from the certification distribution and
nothing else. This becomes decisive in closed loop, where the agent also produces \emph{retries}.

\begin{remark}[Retries are outside the guarantee]\label{rem:retries}
The certification sets contain first proposals. A retry $a_k$, $k\ge1$, differs from them in two ways. Its prompt
$x_k$ contains feedback that a first prompt never contains, which changes the model's internal state. And the
retry is drawn only in episodes whose previous proposal was not executed, which selects harder episodes.
Proposition~\ref{prop:cert} therefore gives no bound on $\rho_A$ for retries. The \emph{first-proposal-only
rule} (R0) replaces \ACC{} by \VAL{} whenever $k\ge1$, so that every unchecked execution is a first proposal, the
population the certificate covers. Unchecked \REJ{} remains allowed on retries, since it cannot execute an
unsafe action.
\end{remark}

\subsection{Cost model and break-even validator cost}\label{sec:costmodel}
Let a policy $p$ (always-validate, or a screen with or without R0) be run on an episode. Its compute is
\begin{equation}\label{eq:cost}
C_p(\cval)=\underbrace{\sum_{k}\bigl(c_{\mathrm{gen},k}+c_{\mathrm{feat},k}\bigr)}_{G_p}\;+\;N_p\,\cval,
\end{equation}
where $c_{\mathrm{gen},k}$ is the generation time of proposal $k$, $c_{\mathrm{feat},k}$ the time of any extra
pass needed to compute features (zero for observables), $N_p$ the number of validator calls and $\cval$ the cost
of one call. In our closed loop the plant is paused while the agent and the validator run, decoding is greedy,
and the validator is deterministic. The sequence of decisions, and with it $G_p$ and $N_p$, is therefore
independent of $\cval$. The expected compute $\bar C_p(\cval)=\bar G_p+\bar N_p\,\cval$ is then exactly linear in
the validator cost. Logged decisions can be re-costed for any validator, faster or slower than the one
measured. Comparing a screen $p$ with always-validate gives the following.
\begin{itemize}
\item If $\bar N_p<\bar N_{\mathrm{AV}}$, the screen saves compute if and only if
\begin{equation}\label{eq:breakeven}
\cval>c^{\ast}_p=\frac{\bar G_p-\bar G_{\mathrm{AV}}}{\bar N_{\mathrm{AV}}-\bar N_p},
\end{equation}
and it saves at any validator cost if $\bar G_p\le\bar G_{\mathrm{AV}}$.
\item The relative saving $1-\bar C_p/\bar C_{\mathrm{AV}}$ increases monotonically with $\cval$ when
$c^{\ast}_p\ge0$ and tends to $1-\bar N_p/\bar N_{\mathrm{AV}}$, the fraction of validator calls avoided, as
$\cval\to\infty$.
\end{itemize}
The numerator of Eq.~\eqref{eq:breakeven} explains when screening is cheap. An \ACC{} of a passing proposal
saves a call and costs nothing, since always-validate would also have executed it. A \REJ{} of a failing
proposal saves a call and also costs nothing in expectation, since always-validate would also have asked for a
new proposal. A \REJ{} of a \emph{passing} proposal, in contrast, forces generations that always-validate would
not have needed, and feature passes add a fixed overhead. Rejection precision therefore governs the break-even
cost, and accept safety governs the risk.

Algorithm~\ref{alg:episode} summarises one closed-loop episode under a screen.

\begin{algorithm}[!htbp]
\caption{One closed-loop episode with a screen (optional R0 rule)}\label{alg:episode}
\begin{algorithmic}[1]
\Require snapshot $s$, agent, probe $f$, thresholds $\tacc,\trej$, validator $V$, budget $K=6$
\State $x_0\gets\mathrm{prompt}(s)$
\For{$k=0,\dots,K-1$}
  \State $(a_k,h_k)\gets\mathrm{agent}(x_k)$ \Comment{proposal and internals}
  \If{$a_k$ repeats an earlier proposal}\ reuse its verdict (no new call) \EndIf
  \State $\hat r\gets f(\psi(x_k,a_k,h_k))$
  \If{$\hat r\le\tacc$ \textbf{and} ($k=0$ \textbf{or} R0 is off)} \State \textbf{execute} $a_k$; \Return \Comment{\ACC}
  \ElsIf{$\hat r\ge\trej$} \State $x_{k+1}\gets\mathrm{prompt}(s,\text{rejection notice})$ \Comment{\REJ: no metrics}
  \Else \State $v\gets V(s,a_k)$ \Comment{\VAL: one simulator call}
    \If{$v=1$} \State \textbf{execute} $a_k$; \Return \EndIf
    \State $x_{k+1}\gets\mathrm{prompt}(s,\text{measured outcome of }a_k)$
  \EndIf
\EndFor
\State hold the current setpoints \Comment{budget exhausted: unresolved}
\end{algorithmic}
\end{algorithm}
